\documentclass[11pt,a4paper]{article}
\usepackage[top=1.6cm, bottom=1.8cm, left=1.8cm, right=1.8cm]{geometry}
\usepackage{amsmath,amssymb,amsfonts}
\usepackage{graphicx}
\usepackage{xcolor}
\usepackage{colortbl}
\usepackage{array}
\usepackage{enumitem}
\usepackage{tikz}
\usepackage{fancyhdr}
\usepackage[T5,T1]{fontenc}
\usepackage[vietnamese]{babel}
\usepackage{newtxtext,newtxmath}

% Định nghĩa màu sắc theo phong cách sách Stewart Calculus
\definecolor{stewartcyan}{RGB}{0, 130, 175}
\definecolor{stewartblue}{RGB}{0, 110, 160}
\definecolor{tableheader}{RGB}{225, 238, 246}
\definecolor{stewartrule}{RGB}{0, 130, 175}

% Thiết lập đầu trang và chân trang
\pagestyle{fancy}
\fancyhf{}
\renewcommand{\headrulewidth}{0pt}
\fancyhead[L]{%
  \fontsize{9pt}{11pt}\selectfont
  \textbf{164}\qquad
  \textbf{\textsf{CHƯƠNG 2}}\quad Giới hạn và Đạo hàm
}

% Định nghĩa biểu tượng màn hình đồ họa CAS (Graphing Calculator)
\newcommand{\casicon}{%
  \begingroup\normalfont
  \begin{tikzpicture}[x=1ex, y=1ex, baseline=-0.2ex]
    \draw[line width=0.6pt, stewartcyan, rounded corners=0.8pt] (0,0) rectangle (1.6, 1.35);
    \draw[line width=0.55pt, stewartcyan] (0.2, 0.3) .. controls (0.55, 1.25) and (0.95, -0.05) .. (1.4, 1.05);
  \end{tikzpicture}%
  \endgroup
}

\setlength{\parindent}{0pt}

\begin{document}

\noindent
\begin{minipage}[t]{0.485\textwidth}

  % Bài 37
  \textbf{37.}\enspace Bảng bên cho biết chiều cao theo thời gian của một cây thông điển hình được trồng để lấy gỗ tại một khu rừng do con người quản lý.
  \par\vspace{2mm}
  \begin{center}
    \setlength{\tabcolsep}{3.5pt}
    \renewcommand{\arraystretch}{1.2}
    \footnotesize
    \begin{tabular}{|>{\cellcolor{tableheader}}l|c|c|c|c|c|c|}
      \hline
      Tuổi cây (năm) & 14 & 21 & 28 & 35 & 42 & 49 \\ \hline
      Chiều cao (feet) & 41 & 54 & 64 & 72 & 78 & 83 \\ \hline
    \end{tabular}
    \par\vspace{1mm}
    {\scriptsize\textit{Nguồn: Arkansas Forestry Commission}}
  \end{center}
  \vspace{1mm}
  Nếu $H(t)$ là chiều cao của cây sau $t$ năm, hãy lập bảng các giá trị ước lượng của $H'$ và phác thảo đồ thị của nó.

  \vspace{3.5mm}
  % Bài 38
  \textbf{38.}\enspace Nhiệt độ nước ảnh hưởng đến tốc độ tăng trưởng của cá hồi suối. Bảng bên cho biết mức tăng khối lượng của cá hồi suối sau 24 ngày ở các mức nhiệt độ nước khác nhau.
  \par\vspace{2mm}
  \begin{center}
    \setlength{\tabcolsep}{4.5pt}
    \renewcommand{\arraystretch}{1.2}
    \footnotesize
    \begin{tabular}{|>{\cellcolor{tableheader}}l|c|c|c|c|c|}
      \hline
      Nhiệt độ ($^\circ\text{C}$) & 15.5 & 17.7 & 20.0 & 22.4 & 24.4 \\ \hline
      Khối lượng tăng (g) & 37.2 & 31.0 & 19.8 & 9.7 & $-9.8$ \\ \hline
    \end{tabular}
  \end{center}
  \vspace{1mm}
  Nếu $W(x)$ là mức tăng khối lượng ở nhiệt độ $x$, hãy lập bảng các giá trị ước lượng của $W'$ và phác thảo đồ thị của nó. Đơn vị của $W'(x)$ là gì?
  \par\vspace{1.5mm}
  {\scriptsize\textit{Nguồn: Phỏng theo J. Chadwick Jr., “Temperature Effects on Growth and Stress Physiology of Brook Trout: Implications for Climate Change Impacts on an Iconic Cold-Water Fish.” Masters Theses. Paper 897. 2012. scholarworks.umass.edu/theses/897.}}

  \vspace{3.5mm}
  % Bài 39
  \textbf{39.}\enspace Gọi $P$ là tỉ lệ phần trăm sản lượng điện của một thành phố được sản xuất bởi các tấm pin năng lượng mặt trời sau $t$ năm kể từ ngày 1 tháng 1 năm 2020.
  \begin{enumerate}[label=\textbf{(\alph*)}, leftmargin=*, topsep=2pt, itemsep=1pt, parsep=0pt]
    \item $dP/dt$ biểu thị điều gì trong ngữ cảnh này?
    \item Hãy giải thích ý nghĩa của phát biểu
    \[
      \left.\frac{dP}{dt}\right|_{t=2} = 3.5
    \]
  \end{enumerate}

  \vspace{2mm}
  % Bài 40
  \textbf{40.}\enspace Giả sử $N$ là số người ở Hoa Kỳ di chuyển bằng ô tô sang một tiểu bang khác để đi nghỉ trong một năm khi giá xăng trung bình là $p$ đô la/gallon. Bạn dự đoán $dN/dp$ là số dương hay số âm? Hãy giải thích.

  \vspace{3.5mm}
  % Nhóm 41-44
  \textbf{\textsf{\color{stewartcyan}41--44}}\enspace Đồ thị của hàm số $f$ được cho như hình vẽ. Hãy chỉ ra, kèm theo lý do, các giá trị mà tại đó $f$ \textit{không} khả vi.
  \par\vspace{2.5mm}

  % Bài 41 và Bài 42
  \noindent
  \begin{minipage}[t]{0.48\linewidth}
    \textbf{41.}\par\vspace{1mm}
    \centering
    \includegraphics[width=\linewidth]{images/p0199_fig5.png}
  \end{minipage}\hfill
  \begin{minipage}[t]{0.48\linewidth}
    \textbf{42.}\par\vspace{1mm}
    \centering
    \includegraphics[width=\linewidth]{images/p0199_fig6.png}
  \end{minipage}

\end{minipage}%
\hfill
\begin{minipage}[t]{0.485\textwidth}

  % Bài 43 và Bài 44
  \noindent
  \begin{minipage}[t]{0.48\linewidth}
    \textbf{43.}\par\vspace{1mm}
    \centering
    \includegraphics[width=\linewidth]{images/p0199_fig1.png}
  \end{minipage}\hfill
  \begin{minipage}[t]{0.48\linewidth}
    \textbf{44.}\par\vspace{1mm}
    \centering
    \includegraphics[width=\linewidth]{images/p0199_fig2.png}
  \end{minipage}

  \par\vspace{2.5mm}
  {\color{stewartrule}\hrule height 0.6pt}
  \vspace{2.5mm}

  % Bài 45
  \noindent\casicon\enspace\textbf{45.}\enspace Vẽ đồ thị của hàm số $f(x) = x + \sqrt{|x|}$. Phóng to liên tiếp nhiều lần, trước hết về phía điểm $(-1, 0)$ và sau đó về phía gốc tọa độ. Có gì khác biệt về hành vi của $f$ ở lân cận của hai điểm này? Bạn rút ra kết luận gì về tính khả vi của $f$?

  \vspace{3mm}
  % Bài 46
  \noindent\casicon\enspace\textbf{46.}\enspace Phóng to về phía các điểm $(1, 0)$, $(0, 1)$ và $(-1, 0)$ trên đồ thị của hàm số $g(x) = (x^2 - 1)^{2/3}$. Bạn nhận thấy điều gì? Hãy giải thích những gì bạn quan sát được dưới góc độ tính khả vi của $g$.

  \vspace{3mm}
  % Nhóm 47-48
  \noindent\textbf{\textsf{\color{stewartcyan}47--48}}\enspace Đồ thị của một hàm số $f$ và đạo hàm $f'$ của nó được cho như hình vẽ. Giá trị nào lớn hơn, $f'(-1)$ hay $f''(1)$?
  \par\vspace{2mm}

  % Bài 47
  \textbf{47.}\par\vspace{1mm}
  \centering
  \includegraphics[width=0.62\linewidth]{images/p0199_fig3.png}
  \par\vspace{3mm}

  % Bài 48
  \raggedright
  \textbf{48.}\par\vspace{1mm}
  \centering
  \includegraphics[width=0.62\linewidth]{images/p0199_fig4.png}

  \vspace{2.5mm}
  {\color{stewartrule}\hrule height 0.6pt}
  \vspace{2.5mm}

  % Bài 49
  \raggedright
  \textbf{49.}\enspace Hình bên cho thấy các đồ thị của $f$, $f'$ và $f''$. Hãy nhận diện từng đường cong và giải thích sự lựa chọn của bạn.
  \par\vspace{2mm}
  \centering
  \includegraphics[width=0.76\linewidth]{images/p0199_fig7.png}

\end{minipage}

\vfill
\begin{center}
  \fontsize{6.5pt}{7.5pt}\selectfont\color{black!70}
  Copyright 2021 Cengage Learning. All Rights Reserved. May not be copied, scanned, or duplicated, in whole or in part. WCN 02-200-203\\
  Chỉ dành cho mục đích giáo dục và học tập cá nhân.
\end{center}

\end{document}